\subsection{The exact question and the specialization comparison}

We retain the entire arithmetic quotient throughout.  If $k$ is a finite
field of characteristic $p$, we put
\[
 \Pi_C=\pi_1(C_{\bar k})^{(p')},\qquad
 1\longrightarrow\Pi_C\longrightarrow E_C^{(p')}
       \longrightarrow G_k\longrightarrow1.
\]
The superscript means a quotient of the \emph{geometric} kernel, not the
maximal prime-to-$p$ quotient of the arithmetic group.  A chosen geometric
point above $P\in C(k)$ gives a representative $s_P$; changing its path
changes it by conjugation in $\Pi_C$.  We call distinct points a collision
only when their section classes agree in this whole extension.

\begin{claim}{AR.1}{Proved; comparison used from R1}
For a smooth proper geometrically connected curve $C/k$ of genus at least
two and $P,Q\in C(k)$, the following are equivalent:
\begin{enumerate}
\item $s_P$ and $s_Q$ are conjugate by $\Pi_C$.
\item For every finite $L/k$ and every connected prime-to-$p$ finite
\'etale Galois cover $D\to C_L$ with constant deck group and a specified
$P'\in D(L)$ above $P$, there is a $Q'\in D(L)$ above $Q$ such that
$[Q'-P']\in J_D(L)[p^\infty]$.
\end{enumerate}
The point $Q'$ in (2) is unique.  One may equivalently allow covers whose
geometric Galois closure has prime-to-$p$ degree.  Thus the covering
condition of cote~149, scan~36, is an exact reformulation of a collision.
Its final implication $P=Q$ is proved below in AR.3.
\end{claim}
\begin{proof}
The forward implication and uniqueness are checked in AC.1.  For a cover
that is only geometrically Galois, use the same fixed-fibre argument:
the open arithmetic subgroup containing the chosen representative of
$s_P(G_L)$ contains the equal representative of $s_Q(G_L)$, which fixes a
point of the fibre above $Q$.  The induced topology on its geometric
group is the full prime-to-$p$ topology.  Indeed a core inside a
prime-to-$p$ Galois closure embeds in a product of prime-to-$p$ groups;
the finite quotient and its extension still have order prime to $p$.
The lifted sections therefore agree in the required full extension.
Prime-to-$p$ Kummer theory gives the asserted torsion class.

For the converse, use $P$ to write $E_C^{(p')}=\Pi_C\rtimes G_k$.
The group $\Pi_C$ is topologically finitely generated by the surface
presentation in SGA~1, XIII, Corollary~2.12 \cite{AR_SGA1}.
It consequently has cofinally many open characteristic subgroups $N$.
The open subgroup $N\rtimes G_k$ defines a geometrically connected
finite \'etale cover $D_N\to C$ with a distinguished rational point
$P_N$ and geometric deck group $\Pi_C/N$.  The continuous arithmetic
action on this finite group becomes trivial over some finite $L/k$.
Hypothesis (2) then gives $Q_N\in D_N(L)$ with $p$-primary difference.
Uniqueness holds over $\bar k$, by AC.1's Hilbert~90 argument.  Every
element of $G_k$ fixes $P_N$ and preserves the torsion condition, so
$Q_N$ is in fact rational over $k$.

For $N'\subset N$ the finite morphism $D_{N'}\to D_N$ sends a point
difference to its image under the norm homomorphism of Jacobians.  The
image of $Q_{N'}$ therefore has the required torsion difference and is
$Q_N$ by uniqueness.  The compatible rational points form a $G_k$-fixed
point in the inverse limit of fibres above $Q$.  This limit is the
prime-to-$p$ \'etale path torsor from $P$ to $Q$.  Its fixed point is
precisely a geometric conjugator between the two sections.  This proves
the converse in the whole profinite group, rather than only at each
uncoordinated finite quotient.  Allowing the additional covers does not
alter the equivalence because the forward implication applies to them.
\end{proof}

\begin{claim}{AR.2}{Proved; the finitely generated field corollary}
Universal absence of collisions for proper curves of genus at least two
over finite fields is equivalent to injectivity of the geometrically
prime-to-$p$ point-section map for every hyperbolic curve over every
field finitely generated over $\mathbf F_p$.  Both assertions hold,
by AR.3 and the comparison proved here.  Independently of AR.3, this
injectivity holds in genus zero and for affine curves of genus one.
For a proper curve over a finite field, $J_C(k)[p]=0$ already suffices.
\end{claim}
\begin{proof}
We verify the comparison in AP-A.8, since a field extension or
specialization cannot silently preserve every rational torsion test.
Over a finitely generated field $K/\mathbf F_p$, the kernel of the
prime-to-$p$ Kummer map of a semiabelian variety is its $p$-primary
torsion.  For its abelian quotient this follows from Lang--N\'eron and
finite generation over a finite constant field.  For a split torus it
follows by taking divisors on a normal proper model: $K^\times$ is a
finite constant group extended by a free abelian group.  Pass to a
finite splitting field for a general torus.  For the extension by an
abelian variety, lifts of finitely many generators leave a free toric
summand, so the same intersection-of-multiples assertion holds.  This
checks the Kummer step of AP-A.1 used here.

The generalized Jacobian of a punctured genus-zero curve is a torus.
Its geometric $p$-primary torsion point group is zero; after splitting
the punctures its Abel map contains a separating cross-ratio
coordinate.  Equality of point sections is therefore impossible for
distinct points.  For an affine genus-one curve $E\setminus B$, choose
the alleged colliding point $P$ as origin.  Its difference from $Q$
has order $p^a$.  Choose $m$ prime to $p$ with $m^2>p^a$ and use the
finite \'etale cover
\[
 E\setminus[m]^{-1}B\ \xrightarrow{[m]}\ E\setminus B.
\]
After a finite separable extension defining its deck group and boundary,
a collision lifts to $0,Q'$.  The generalized-Jacobian difference has
order $p^a$: its elliptic projection has that order and its toric kernel
has no $p$-primary geometric points.  There is a rational function with
divisor $p^a(Q'-0)$ and value $1$ at every boundary point.  Subtracting
$1$ gives at least $m^2$ distinct zeros to a function of degree $p^a$,
a contradiction.  These arguments include $p=2,3$.

For genus at least two, passage from an affine curve to its smooth
proper compactification induces a continuous quotient of arithmetic
extensions over $G_K$ and maps geometric conjugators to geometric
conjugators.  It remains to specialize a proper collision.  Spread the
smooth proper curve and its two distinct sections over a smooth
connected finite-type scheme $S/\mathbf F_q$ and shrink so that the
sections are disjoint.  At a closed point $s$, put
$R=\mathcal O_{S,s}^{h}$, $L=\operatorname{Frac}(R)$ and
$k_s=\kappa(s)$.  The map $G_L\twoheadrightarrow\pi_1(R)=G_{k_s}$ is
surjective because a connected finite \'etale cover of the normal
integral henselian local scheme remains connected generically.
SGA~1, XIII, Proposition~4.3, Examples~4.4 and 4.5.3 give the
specialization comparison for the smooth proper family with section;
its geometric prime-to-$p$ group is constant under this comparison.
Thus there is a continuous homomorphism
\[
 E_{C_L}^{(p')}\longrightarrow E_{C_s}^{(p')}
\]
which is an isomorphism on geometric kernels, induces the displayed
surjection on arithmetic quotients, and sends the section of each
specialized rational point to its section after that quotient map.
An equality of section representatives upstairs consequently gives an
equality downstairs.  The sections remain distinct as points.  The
residue field is finite, which proves the required reduction.  The
reverse universal implication includes finite fields among finitely
generated fields.  Finally, for a finite field the Kummer kernel is
$\bigcap_{(n,p)=1}nJ_C(k)=J_C(k)[p^\infty]$; the Abel embedding
settles the last assertion when this group is zero.
\end{proof}

\begin{claim}{AR.3}{Proved; universal injectivity of proper point sections}
\label{clm:ARuniversalpointsections}
Every smooth proper geometrically connected curve of genus at least two
over every finite field has an injective geometrically prime-to-$p$
point-section map.  Equivalently, the last implication in AR.1 always
holds, or the universal conclusion of AR.2 holds.
The arithmetic quotient is the whole $G_k$, and conjugation is by the
whole geometric prime-to-$p$ group.  In particular the assertion includes
characteristics two and three and every finite constant extension.
\end{claim}
\begin{proof}
We assemble the three group-level comparisons proved in AG.40--AG.42.
Suppose $P,Q\in C(k)$ have geometrically conjugate sections, and choose
any prime $r\ne p$.  AG.41 proves exactness for the configuration fiber
after retaining its entire proper geometric prime-to-$p$ quotient and
replacing only the cuspidal kernel by its maximal pro-$r$ quotient.
Its arithmetic comparison retains $G_k$ and identifies each punctured
curve group with the actual inverse image of its point section under
the second projection of the configuration group.

AG.42 identifies the puncture decomposition subgroup with the
intersection with the global diagonal boundary group.  A geometric
conjugator between $s_P(G_k)$ and $s_Q(G_k)$ lifts to that boundary
group.  Conjugation by this lift, followed by a geometric inner
correction in the target fiber, gives a continuous isomorphism
between the maximal cuspidally pro-$r$ groups of $C-P$ and $C-Q$.
It is over the identity of the full proper arithmetic quotient and
carries the actual cusp decomposition group to a conjugate of the
actual target cusp group.  AG.40 proves that such an isomorphism
forces $P=Q$.  This contradicts the existence of a collision.

For clarity, the implication in AG.40 uses the entire family of finite
covers visible in this cuspidal quotient.  Their full same-prime
compact cohomology and Frobenius are preserved.  Their zeta functions
recognize noncuspidal point sections by a compatible inverse-limit
argument; the degree-one cuspidal central extension then turns this
recognition into a Kummer evaluation contradiction.  No finite-level
test is substituted for whole-group conjugacy.

The full prime-to-$p$ configuration fiber map is itself noninjective,
as proved in AG.39.  The fixed-$r$ cuspidal exactness used here is a
different statement with its own finite-quotient proof.  The proof
does not use the disputed full configuration assertion in the
cuspidalization literature.  Together with AR.1 and AR.2, it supports
the injectivity contemplated in cote~149, scans~35--36, including
the question ``Il se pourrait qu'on puisse en d\'eduire $P=Q$''.
\end{proof}
